% Propietario reproducido por unidad demostrativa; véase PROPIETARIOS_CONTINUO_UNIVERSO.json.
% Origen: XI ES CUMPLIMIENTO_EDITORIAL_20260928, sections/ch_cierre_cardinal.tex:26--403.
% Cuerpos y pruebas conservados; sólo namespace, rutas y remisiones contextuales declaradas.
\subsection{Extension of states and transfinite semantic extension}

Let

\[
 \Sigma_{\rm HMT}
 =\bigl(\dot R_{\rm APP},\dot R_{\rm TRIT},\dot R_{\rm TPK},
 \dot R_{U},\dot R_{\Gamma_9},\dot R_{\rm Ext},\dot R_\rho,
 \dot R_{\pi},\dot R_{\varphi},\dot R_e,\dot R_{\alpha},
 \dot R_{\rm ar},\dot R_{\rm exc}\bigr)
\tag{33}\label{hmtfg:base:eq:ch-signatura}
\]

be the countable \textbf{relational} signature whose effective code is \(h\).
Each \(\dot R_S\) is the graph of the operator \(S\) on the codes where it is
defined; no external function is assumed to be total within each level. The
graphs \(\dot R_{\pi},\dot R_{\varphi},\dot R_e,\dot R_{\alpha}\) are the
output readers constructed by coinductive generation and dodecaphase closure;
they do not take the conventional values of the constants as arguments. Their
presence in \(\Sigma_{\rm HMT}\) ensures that the closure preserves the full
state: it simultaneously integrates the routes that express value and the
incidence that extends toward the exceptional realization.

At each horizon, the state code records the APP--TRIT--TPK coordinates, the
route, sheet, carry, boundary, and memory; at its representation stages it also
records \(P_N,\Phi_N,E_N,A_N\), where \(A_N\) is the prefix of
\(\alpha_{\rm HMT}\) produced by the same state. These coordinates do not
select the branch from outside: they are part of the history being encoded.
Once a primitive-recursive numbering of the finite coordinates of a state is
fixed, every inverse history
\(x_\infty=(x_N)_{N\geq1}\in X_\infty^{\rm enr}\) determines the code
\[
 \operatorname{code}(x_\infty)
 :=\{\langle N,j,c\rangle:(x_N)_j\text{ has code }c\}\subseteq\omega.
\]
The coherence equations
\(\rho_{N+1,N}(x_{N+1})=x_N\) make it possible to decode the history uniquely
from that real. We write \(m_\infty\) for the code of the realized history and
\(m_\infty(n)\) for its characteristic function under the fixed numbering.
The operator code \(h\) and this complete ledger are combined, by means of a
primitive-recursive pairing function, into the real

\[
 u_{h,m}:=
 \{\langle0,\langle n,h(n)\rangle\rangle:n<\omega\}
 \cup
 \{\langle1,\langle n,m_\infty(n)\rangle\rangle:n<\omega\}.
\tag{33a}\label{hmtfg:base:eq:ch-codigo-conjunto}
\]

The two primitive projections recover \(h\) and \(m_\infty\) from
\(u_{h,m}\).

The first extension is the metric and operator extension of the finite states.
For \(N\ge1\), let \(\widehat{\mathcal X}^{\rm enr,raw}_{6N}\) be the set of
raw enriched states of length \(6N\), and define the surviving subsystem
set-theoretically by
\[
\begin{aligned}
 X_N
 &:=\widehat{\mathcal X}^{\rm enr,surv}_{6N}\\
 &:=\Bigl\{x\in\widehat{\mathcal X}^{\rm enr,raw}_{6N}:\\[-.2ex]
 &\qquad\forall M>N\ \exists z\in
   \widehat{\mathcal X}^{\rm enr,raw}_{6M}\;\text{ such that}\\[-.2ex]
 &\qquad\operatorname{tr}_{6M,6N}(z)=x\ \wedge\
   z\text{ satisfies every intermediate TPK transition}\Bigr\}.
\end{aligned}
\]
and set
\[
 \rho_{N+1,N}:=\operatorname{tr}_{6N+6,6N}\!\upharpoonright X_{N+1},
\]
and
\[
 \operatorname{Ext}_N(x):=
 \{y\in X_{N+1}:(x,y)\in
 \operatorname{Ext}_{6N+6,6N}\}.
\]
The root level is included without any implicit exception:
\[
 X_0:=\{\ast\},\qquad
 \rho_{1,0}:X_1\to X_0\text{ is constant},\qquad
 \operatorname{Ext}_0(\ast):=X_1.
\]
The superscript \(\mathrm{surv}\) denotes the pruned tree of states possessing
compatible TPK descendants at every finite horizon. The TPK extension law
established in the preceding sections proves that the emitted seeds belong to
this tree; finite branching and Kőnig's lemma turn compatibility at every
horizon into an infinite branch. In particular,

\[
 \forall N\ge0\;\forall x\in X_N\quad
 \varnothing\ne\operatorname{Ext}_N(x)
 \subseteq\rho_{N+1,N}^{-1}(x).
\tag{34}\label{hmtfg:base:eq:ch-extension-no-vacia}
\]

By the same finitely branching tree lemma, every surviving seed has a
compatible extension

\[
 x_\infty=(x_N)_{N\ge1}\in
 X_\infty^{\rm enr}:=\varprojlim_N(X_N,\rho_{N+1,N}).
\tag{35}\label{hmtfg:base:eq:ch-limite-inverso}
\]

The second extension acts \textbf{afterward} on the complete code of that
realized history.
It is not inferred from the mere existence of the finite operator
\(\operatorname{Ext}_N\): at the set-theoretic scale, it formalizes the HMT principle that
only those objects possessing a genealogical derivation
from the state and its rules belong to the full domain. The transfinite semantic
extension is:

\[
 \begin{aligned}
  \mathfrak G_0(h,m_\infty)
  &:=\operatorname{TC}\bigl(\{\omega,u_{h,m},h,m_\infty\}\bigr),\\
  \mathfrak G_{\beta+1}(h,m_\infty)
  &:=\operatorname{Def}_{\Sigma_{\rm HMT}}
       \bigl(\mathfrak G_\beta(h,m_\infty)\bigr),\\
  \mathfrak G_\lambda(h,m_\infty)
  &:=\bigcup_{\beta<\lambda}\mathfrak G_\beta(h,m_\infty)
       \quad(\lambda\text{ limit}),\\
  \mathfrak G_{\rm HMT}(h,m_\infty)
  &:=\bigcup_{\beta\in\operatorname{Ord}}
       \mathfrak G_\beta(h,m_\infty).
 \end{aligned}
\tag{36}\label{hmtfg:base:eq:ch-clausura-transfinita}
\]

The Continuum Hypothesis and a target bijection are not introduced among the
conditions defining the successor operator; the operator is given
prospectively by

\[
 \operatorname{Def}_{\Sigma_{\rm HMT}}(M)
 :=M\cup
 \Bigl\{
   \{x\in M:\langle M,\in,
       (\dot R_S\cap M^{k_S})_{S\in\Sigma_{\rm HMT}}\rangle
       \models\varphi(x,\bar p)\}:
   \ulcorner\varphi\urcorner\in\omega,
   \ \bar p\in M^{<\omega}
 \Bigr\}.
\tag{37}\label{hmtfg:base:eq:ch-operador-def}
\]

That is, it simultaneously expresses all subsets definable by a finite formula
of the HMT relational language and parameters already produced by the
genealogy. It consults neither a target real value, a conventional constant, a
target cardinal, nor an external model. Formula~\eqref{hmtfg:base:eq:ch-operador-def}
defines the transfinite closure over the previously constructed seeds,
operators, enriched state, ledger, and double projection.

\begin{lemma}[Joint coding of the HMT program and the realized ledger]
\label{hmtfg:base:thm:ch-6}
The pair formed by the operator code \(h\) and the ledger \(m_\infty\) is
encoded by a single real \(u_{h,m}\subseteq\omega\), and every formula of
\(\Sigma_{\rm HMT}\) is uniformly translated into a formula of the language
\(\{\in,u_{h,m}\}\).

\end{lemma}
\begin{proof}
Both \(h\) and \(m_\infty\) are sequences of natural numbers: \(h\) encodes
the signature, the finite tables, and the operator programs; \(m_\infty\)
encodes one realized history, not the complete history space.
Formula~\eqref{hmtfg:base:eq:ch-codigo-conjunto} and its two decoders prove the first
assertion. Each symbol in~\eqref{hmtfg:base:eq:ch-signatura} denotes its \textbf{graph
relation encoded in \(h\)}, never the external graph of all its possible
evaluations. APP, TRIT, \(U_t\), \(\Gamma_9\), the restrictions \(\rho\), and
the finite extensions have primitive-recursive graphs on integer codes and,
therefore, \(\Delta_0\) formulas absolute between transitive levels. The two
representations are expressed as graph relations through the formulas

\[
 \begin{aligned}
 \dot R_{\rm ar}(x,y)
 &\;\Longleftrightarrow\;
 \forall k\in\omega\;\exists N\;\forall n\ge N\quad
 |\mu_n(x)-y|<2^{-k},\\
 \dot R_{\rm exc}(x,c,z)
 &\;\Longleftrightarrow\;
 \forall N\in\omega\quad z\!\upharpoonright N
   =E_N(x\!\upharpoonright N,c\!\upharpoonright N),
 \end{aligned}
\tag{37c$_0$}\label{hmtfg:base:eq:ch-grafos-publicacion}
\]

where \(\mu_n\) and \(E_N\) are the rational/finite functions whose codes
occur in \(h\). Every atomic formula containing an HMT symbol is inductively
replaced by its graph formula with parameter \(u_{h,m}\). Connectives and
quantifiers are preserved. This induction on syntactic complexity produces a
translation \(\varphi\mapsto\varphi^*\) such that, for every transitive level
\(M\) of~\eqref{hmtfg:base:eq:ch-clausura-transfinita} containing the parameters,

\[
 \langle M,\in,(\dot R_S\cap M^{k_S})_S\rangle
   \models\varphi(\bar a)
 \quad\Longleftrightarrow\quad
 M\models\varphi^*(\bar a;u_{h,m}).
\tag{37c}\label{hmtfg:base:eq:ch-eliminacion-uniforme}
\]

Thus, \eqref{hmtfg:base:eq:ch-operador-def} cannot consult an oracle, a target decimal
table, or an uncoded relation that introduces uncountably many objects in a
single step.
\end{proof}

HMT semantics is now defined \textbf{independently of identifying it with
\(\mathfrak G\)}. Let \(\mathcal T_0^{\rm HMT}\) be the set of canonical names
of the elements of
\(\mathfrak G_0(h,m_\infty)=
\operatorname{TC}(\{\omega,u_{h,m},h,m_\infty\})\). In particular, the
initial level contains the realized ledger \(m_\infty\), not every history
in \(X_\infty^{\rm enr}\) as primitive parameters. The admissible terms are
the well-founded trees obtained by

\[
 \begin{aligned}
 \mathcal T^{\rm HMT}_{\beta+1}
 &:={\mathcal T}^{\rm HMT}_\beta\cup
 \{\langle\beta,\ulcorner\varphi\urcorner,
       t_0,\ldots,t_{k-1}\rangle:
       t_i\in\mathcal T^{\rm HMT}_\beta\},\\
 \mathcal T^{\rm HMT}_\lambda&:=
       \bigcup_{\beta<\lambda}\mathcal T^{\rm HMT}_\beta,
 \qquad
 \mathcal T^{\rm HMT}:=
       \bigcup_{\beta\in\operatorname{Ord}}\mathcal T^{\rm HMT}_\beta.
 \end{aligned}
\tag{37a}\label{hmtfg:base:eq:ch-terminos}
\]

The value of a successor term is the subset of the preceding level defined by
\(\varphi\) in the induced relational structure, with the values of
\(t_0,\ldots,t_{k-1}\) as parameters. Before any set-theoretic comparison,
define

\[
 \operatorname{Ob}_{\rm sem}^{\rm HMT}(h,m_\infty)
 :=\{\llbracket t\rrbracket:t\in\mathcal T^{\rm HMT}\}.
\tag{37b}\label{hmtfg:base:eq:ch-objetos-semanticos}
\]

This is the rule of \textbf{generative exhaustiveness} adopted by HMT: every object in the entire domain must have one of those derivation trees and every admissible tree expresses an object. It is not a consequence of finite branching. The continuum hypothesis is not introduced into this rule; it is obtained through the proof developed in the article. Neither an enumeration of the reals nor a target bijection is introduced as data.

\begin{theorem}[Equivalence between term semantics and the genealogical hierarchy]
\label{hmtfg:base:thm:ch-6a}
The evaluation of the preceding terms satisfies

\[
 \operatorname{Ob}_{\rm sem}^{\rm HMT}(h,m_\infty)
 =\mathfrak G_{\rm HMT}(h,m_\infty).
\tag{37d}\label{hmtfg:base:eq:ch-equivalencia-semantica}
\]

\end{theorem}
\begin{proof}
By transfinite induction. At the initial level, the two classes coincide. If
every value of a term in \(\mathcal T^{\rm HMT}_\beta\) belongs to
\(\mathfrak G_\beta\), the semantics of a successor term is one of the subsets
in~\eqref{hmtfg:base:eq:ch-operador-def}. Conversely, every subset in
~\eqref{hmtfg:base:eq:ch-operador-def} is given by a formula code and a finite tuple of
parameters; by the induction hypothesis those parameters possess terms, and
~\eqref{hmtfg:base:eq:ch-terminos} forms the term that denotes it. At a limit, both
constructions take the same union.
\end{proof}

\Needspace{9\baselineskip}
The state encoder is compatible with truncations and sends every TPK extension
to a successor term. Let
\(\mathcal T_N^{\rm st}\subset\mathcal T^{\rm HMT}\) be the subclass of terms
encoding complete states of depth \(N\), and let
\(\tau_{N+1,N}:\mathcal T_{N+1}^{\rm st}\to\mathcal T_N^{\rm st}\) delete the
last block together with the corresponding truncation of the genealogical
ledger, carry, and memory. If \(J_N:X_N\to\mathcal T_N^{\rm st}\) is the
complete encoder, compatibility with \(\operatorname{Ext}_N\) proves the
well-typed identity

\[
 \forall x\in X_N\ \forall y\in\operatorname{Ext}_N(x)\qquad
 \tau_{N+1,N}\bigl(J_{N+1}(y)\bigr)=J_N(x).
\tag{37e$_0$}\label{hmtfg:base:eq:ch-naturalidad-codificador}
\]

Define within the semantics
\[
 \operatorname{Succ}^{\rm TPK}_{N+1}(J_N(x)):=
 \{t\in\mathcal T_{N+1}^{\rm st}:\exists y\in X_{N+1}\,
   (y\in\operatorname{Ext}_N(x)\ \wedge\ t=J_{N+1}(y))\}.
\]
Choose \(\beta\) such that \(\mathfrak G_\beta\) contains \(x\), \(X_N\),
\(X_{N+1}\), the finite graphs of
\(\operatorname{Ext}_N,J_N,J_{N+1}\), and their codes in \(h\). Then the
primitive-recursive formula for the graph of \(\operatorname{Ext}_N\) gives

\[
 J_{N+1}[\operatorname{Ext}_N(x)]
 =\operatorname{Succ}^{\rm TPK}_{N+1}(J_N(x)),
 \qquad
 J_{N+1}[\operatorname{Ext}_N(x)]
 \in\mathfrak G_{\beta+1},
\tag{37e}\label{hmtfg:base:eq:ch-sucesores-semanticos}
\]

and both sides are formed by the graph formula of the same TPK extension.
Equations~\eqref{hmtfg:base:eq:ch-naturalidad-codificador}--\eqref{hmtfg:base:eq:ch-sucesores-semanticos}
are the explicit semantic link: the subclass of successor terms is exactly the
image of the admissible extensions, and its truncation commutes. The complete
fiber of \(\rho\) is not identified with \(\operatorname{Ext}_N\), nor is
\(\operatorname{Def}\) identified with TPK merely by nomenclature.

\begin{lemma}[Transitivity, axioms, and well-ordering of the genealogical model]
\label{hmtfg:base:thm:ch-6b}
\(\mathfrak G_{\rm HMT}(h,m_\infty)\) is a transitive class, contains all
ordinals, satisfies Zermelo--Fraenkel set theory (ZF), and possesses a
definable global well-order; it therefore satisfies Choice and all axioms used
in the subsequent cardinal comparison.

\end{lemma}
\begin{proof}
Abbreviate \(B=\mathfrak G_0\) and \(u=u_{h,m}\). Induction on
~\eqref{hmtfg:base:eq:ch-clausura-transfinita} proves simultaneously that every level is
transitive and that \(\mathfrak G_\alpha\in\mathfrak G_{\alpha+1}\); if all
ordinals smaller than \(\alpha\) are present, their set is definable at the
corresponding level and expresses \(\alpha\). Hence the union contains all
ordinals. Extensionality and Foundation are inherited from the ambient
universe; Empty Set and Infinity belong to \(B\). Once a level containing the
parameters is chosen, the formulas defining \(\{a,b\}\), \(\bigcup a\), and
subsets by Separation express them at a successor level.

For Collection--Replacement, let \(F\) be a functional relation definable in
\(\mathfrak G\) on the set \(a\). The reflection scheme relative to the
definable hierarchy
\((\mathfrak G_\xi,\in,u_{h,m})_{\xi\in\mathrm{Ord}}\), applied in the ambient
metatheory to the finite list of formulas defining \(F\), produces a level
\(\mathfrak G_\theta\) containing \(a\) and the parameters and correct for
those formulas. The ranks of the witnesses in \(F[a]\) are then bounded by
\(\theta\); Separation over \(\mathfrak G_\theta\) forms the image in
\(\mathfrak G_{\theta+1}\). For Power Set, the internal set to be constructed
is not presupposed. Given \(a\in\mathfrak G_\gamma\), consider in the ambient
metatheory the already existing set \(\mathcal P^V(a)\). For every
\(b\in\mathcal P^V(a)\cap\mathfrak G\), let \(r(b)\) be its first
genealogical stage. Since \(\mathfrak G\) is a definable class, Separation in
the ambient metatheory first forms the set
\(\mathcal P^V(a)\cap\mathfrak G\). Replacement \textbf{in the ambient
metatheory}, applied to that set, bounds the ordinals \(r(b)\) by some
\(\theta\). Consequently,
\[
 \mathcal P^{\mathfrak G}(a)
 =\{b\in\mathfrak G_\theta:b\subseteq a\},
\]
and the right-hand side is expressed by Separation in
\(\mathfrak G_{\theta+1}\). It does not introduce into \(\mathfrak G\) the
subsets external to the semantics~\eqref{hmtfg:base:eq:ch-objetos-semanticos}. This is the
noncircular rank-bounding argument.

Finally, each object is assigned its first level of birth and its least
\textbf{derivation name}: a well-founded, finitely branching derivation tree
(and, by~\eqref{hmtfg:base:eq:ch-terminos}, finite for each individual name), labeled by
a level ordinal, a formula code, and the tuple of names of its parameters.
These names are ordered by transfinite recursion: first by the ordinal label,
then by the formula code, and then lexicographically by the already ordered
parameter names. The least name of each value defines a class relation
\(<_{\mathfrak G}\) that compares all objects. It is a definable global
well-order and, by taking the least element of every nonempty set, proves
Choice.
\end{proof}

% Propietario reproducido por unidad demostrativa; véase PROPIETARIOS_CONTINUO_UNIVERSO.json.
% Origen: XI ES CUMPLIMIENTO_EDITORIAL_20260928, sections/ch_reflexion_testigos.tex:1--89.
% Cuerpos y pruebas conservados; sólo namespace, rutas y remisiones contextuales declaradas.
\paragraph{Reflection by witness closure and construction of the image.}
The Collection--Replacement step of Lemma~\ref{hmtfg:base:thm:ch-6b} admits the following
explicit proof. We retain the hierarchy already generated from the operator
code and the realized record, and abbreviate
\(\mathfrak G=\mathfrak G_{\rm HMT}(h,m_\infty)\).
The argument uses Separation and Replacement in the ambient metatheory,
without presupposing these schemes in \(\mathfrak G\).

Let \(\Phi\) be a finite collection of formulas. First translate their
universal quantifiers using negation and existential quantification, and
close the resulting collection under subformulas. For every initial level
there is a later limit ordinal \(\theta\) such that
\[
 \mathfrak G_\theta\models\varphi(\bar a)
 \quad\Longleftrightarrow\quad
 \mathfrak G\models\varphi(\bar a)
 \qquad
 (\varphi\in\Phi,\ \bar a\in\mathfrak G_\theta^{<\omega}).
\]
For each fixed formula, satisfaction in the class is interpreted by
relativizing its quantifiers to the definable class \(\mathfrak G\);
no uniform truth predicate is introduced.

\begin{proof}
Fix \(\alpha\). For each existential subformula
\(\exists y\,\psi(y,\bar x)\) of \(\Phi\), the tuples
\(\bar a\in\mathfrak G_\alpha^{<\omega}\) of the corresponding arity
that have a witness in \(\mathfrak G\) form a set by Separation in the
metatheory. Assign to each tuple the least ordinal \(\beta\) for which
some \(y\in\mathfrak G_\beta\) satisfies
\(\mathfrak G\models\psi(y,\bar a)\). This ordinal exists because
\(\mathfrak G\) is the union of its levels. Replacement in the metatheory
bounds these first stages of witnesses. Since \(\Phi\) is finite, a
single bound \(f(\alpha)>\alpha\), valid for all its existential
subformulas, can be chosen definably. Thus every tuple under consideration
has \emph{at least one witness} in \(\mathfrak G_{f(\alpha)}\).

Take \(\alpha_0\) large enough to contain the initial parameters and set
\[
 \alpha_{n+1}=f(\alpha_n),\qquad
 \theta=\sup_{n<\omega}\alpha_n.
\]
By continuity of the hierarchy, every finite tuple from
\(\mathfrak G_\theta\) belongs to some \(\mathfrak G_{\alpha_n}\).
If an existential formula of \(\Phi\) with that tuple is true in
\(\mathfrak G\), it has a witness in
\(\mathfrak G_{\alpha_{n+1}}\subseteq\mathfrak G_\theta\).
Induction on formulas now proves the equivalence: atomic formulas are
preserved in the induced structure; connectives preserve the equivalence;
the existential step uses that witness in one direction and the same
witness, viewed in \(\mathfrak G\), in the other. The prior translation
of universal quantifiers also includes the existential formulas expressing
their possible counterexamples. This proves the required finite reflection.
\end{proof}

Now let \(F(x,y,\bar p)\) be functional on \(a\) in \(\mathfrak G\).
Apply the construction to the formulas expressing \(F\), the existence
of its witnesses, and its image, with \(a,\bar p\in\mathfrak G_\theta\).
Transitivity gives \(a\subseteq\mathfrak G_\theta\). Every
\(x\in a\) has its value within that level, and reflection identifies
the exact image as
\[
 b=\{y\in\mathfrak G_\theta:
       \mathfrak G_\theta\models
       \exists x\in a\ F(x,y,\bar p)\}.
\]
This set is definable over \(\mathfrak G_\theta\), with parameters from
the same level. By the successor operation of the hierarchy,
\(b\in\operatorname{Def}(\mathfrak G_\theta)
\subseteq\mathfrak G_{\theta+1}\). This establishes Replacement in
\(\mathfrak G\). If the existence of witnesses is retained and uniqueness
is omitted, the same construction provides a set of witnesses for Collection.

\paragraph{Arity of names and internal selection.}
In the well-order of Lemma~\ref{hmtfg:base:thm:ch-6b}, the enumeration of the base
assigns to each element its least natural-number code. At each successor
stage, formula codes are ordered first; once a code is fixed, the arity
of its parameter tuple is fixed. Lexicographic order is then applied to a
\emph{finite product of fixed arity} of the previously constructed
well-orders. The least name represents each new value, and limit levels
retain the order of first appearance. This specifies the ordering of
tuples without combining different arities into a single lexicographic
order. Selecting the least element of every nonempty set, together with
Replacement already established, forms internally the choice functions
used subsequently. This construction of names precedes the cardinal
injections and receives no bijection of the continuum as data.

